\begin{abstract}
\noindent Every causal conclusion drawn from observational data is conditional on a causal graph, and in most applied work the graph is asserted rather than tested. This paper identifies a class of administrative datasets for which the testing problem is tractable: compliance-gated datasets, in which units report outcomes if and only if observable characteristics cross a regulatory threshold, so that the institutional documents defining the reporting obligation also encode the process that generated the data. We develop a three-stage pipeline, the Tested Data-Generating Process (TDGP), that locates the compliance gate in the dataset's missingness structure, formulates a candidate causal graph from that structure and the variable descriptions alone, and then tests the candidate against the constitutive documents, which are quarantined until the test. We demonstrate the pipeline on the Chicago Energy Benchmarking dataset, 28,329 building-year records (2014--2023) gated at 50,000 square feet, where the candidate survives the documentary test and the recorded data reproduce the EPA greenhouse-gas accounting identity to $R^2 = 0.9999$. A 2$\times$3 factorial experiment crossing causal correction with functional form then shows what the tested structure licenses: an inverse-probability-weighted hierarchical Bayesian model and IPW XGBoost converge to within 0.4 percentage points of median absolute percentage error (21.4\% against 21.0\%), evidence that structural and algorithmic models estimate the same causal quantity once the data-generating process is tested, and a resolution of Breiman's two-cultures tension under that condition. The tested structure further supports deployment prediction for the 5,510 non-compliant records excluded from all training and counterfactual policy queries that the gated sample alone cannot answer.
\par\bigskip\noindent\textbf{Keywords:} \textit{causal inference; directed acyclic graphs; selection bias; missing data; inverse probability weighting; energy benchmarking}
\end{abstract}
